\chapter{State-space models, selective updates and scans}

A recurrence appears to demand one update after another. Before computing
$h_4$, surely we must compute $h_3$, then $h_2$, then $h_1$?
That is true of one way to execute the recurrence. It is not always true
of the underlying computation. Sometimes a whole interval can be
represented by a small summary that composes with the summary of another
interval.

This chapter builds that summary from first principles. We begin with a
scalar state, derive a causal convolution, introduce input-dependent
updates, and implement a tree-based prefix computation. We then test
forward values, gradients, padding, and chunk boundaries. The result is
a reference implementation of diagonal affine scans, not a reproduction
of Mamba, S4, S5, or RWKV, and not a trained genomic model.

We assume Chapter 6's distinction between a numerical state and its
gradient history, and Chapter 5's sequence representation contract.
The new mathematics is exponential decay and associative composition.
The important conceptual distinction is between \emph{linear in the
previous state} and \emph{linear in the entire input sequence}.

\section{Start with a state whose arithmetic is visible}

Consider the scalar update
\begin{equation}
 h_t=a_t h_{t-1}+b_t.
 \label{eq:affine}
\end{equation}
The pair $(a_t,b_t)$ describes a function: multiply the incoming state
by $a_t$, then add $b_t$. The additive term may already include an
input projection. For several independent state coordinates, multiplication
is elementwise and $a_t,b_t,h_t\in\mathbb R^N$.

Here is a four-step example with $h_0=2$. Predict the final state before
reading the last column. The second step is an identity update; it
should not erase or advance the stored value.
\begin{center}\input{\Generated/results-table.tex}\end{center}
The states are $2,2,2.5,1$. Columns $P_t,Q_t$ summarize the entire prefix:
$h_t=P_th_0+Q_t$. We will derive those columns rather than accepting
them as another execution trace.

\Listing{sequential}

The reference uses shape $[T,N]$ for one sequence, with initial state
$[N]$. A batch adds an independent leading batch dimension; no sequence
should share another row's state. The implementation validates shapes,
finite inputs, dtype, and device. An empty sequence returns an empty
state sequence. Its incoming state remains the caller's carry value.

\Plate{01-state-update}{One diagonal affine state update. Retention and input-derived write values meet at an addition; a separate readout maps state to output. Shapes and dependencies are explicit. These are computational states, not biological concentrations or genes.}{fig:update}

A full sequence model needs more than this recurrence: input representation,
mixing across channels, nonlinear transformations, readout, loss, and a
training procedure. We deliberately isolate one operation so its numerical
meaning can be checked before adding those layers.

\section{Derive a stable discrete update from continuous time}

For one coordinate, let an internal continuous state satisfy
\begin{equation}
 \frac{dh}{ds}=-\lambda h+u,\qquad \lambda>0.
 \label{eq:continuous}
\end{equation}
Hold $u$ constant over an interval of length $\Delta$.
Multiplication by the integrating factor $e^{\lambda s}$ gives
$\frac{d}{ds}(e^{\lambda s}h)=e^{\lambda s}u$.
Integrating over that interval yields
\begin{equation}
 h_{\rm new}=e^{-\lambda\Delta}h_{\rm old}
       +\frac{1-e^{-\lambda\Delta}}{\lambda}u.
 \label{eq:zoh}
\end{equation}
This is \Term{zero-order hold} discretization: constant input on the
interval, exact integration of this linear system on that interval.
It is not the same as a forward Euler step,
$h_{\rm new}=(1-\lambda\Delta)h_{\rm old}+\Delta u$.

For $\lambda\Delta=3$, Euler's homogeneous multiplier is $-2$:
successive steps alternate sign and grow. The exact multiplier is
$e^{-3}\approx0.0498$, a contraction. The continuous system was stable;
the wrong step approximation introduced the instability.

\Listing{zoh}

The function uses \texttt{expm1} to compute the small difference
$e^z-1$ accurately. For very small positive $\Delta$, the write gain
approaches $\Delta$. At $\Delta=0$, the pair is exactly $(1,0)$:
no internal time passes and no input is written. We require $\lambda>0$;
the $\lambda=0$ limiting case is an integrator and needs an explicit
branch rather than division by zero.

\Plate{02-timescales}{Computed homogeneous retention for two positive decay rates at unit intervals. Each curve starts from the same state and has zero drive. Different coordinates can represent different timescales; long retention alone does not establish useful memory or biological understanding.}{fig:timescales}

For a learned sequence model, $s$ and $\Delta$ can be internal computational
coordinates. They need not represent physical seconds or genomic base pairs.
Changing tokenization changes what one sequence step means. A mathematical
decay constant should not be described as a measured biochemical lifetime.

In a vector system the familiar notation is
$\dot h=Ah+Bx$, $y=Ch$, with a discretized update
$h_t=\bar A h_{t-1}+\bar Bx_t$. Our derivation covers a diagonal negative
$A$ and a scalar drive per coordinate. General matrix exponentials and
structured parameterizations require additional numerical machinery.
The S4 paper introduces a structured sequence-model construction, not
merely a new name for any exponentially decaying scalar \cite{s4}.

\section{Expand a fixed recurrence into a convolution}

Freeze the coefficients: $a_t=a$, $b_t=gx_t$, and let $y_t=ch_t$.
For zero initial state,
\begin{align}
 h_1&=gx_1,\\
 h_2&=agx_1+gx_2,\\
 h_3&=a^2gx_1+agx_2+gx_3.
\end{align}
In general,
\begin{equation}
 y_t=\sum_{j=1}^t cga^{t-j}x_j,\qquad K_\ell=cga^\ell.
 \label{eq:convolution}
\end{equation}
The coefficient depends on lag $t-j$, not the absolute positions.
This is a causal \Term{linear time-invariant} convolution. It gives
another way to compute the same transformation when the entire input
is available.

\Listing{convolution_form}

This direct implementation intentionally costs $O(T^2)$; it is an
independent, readable oracle rather than a fast convolution kernel.
For $a=0.5$, $g=c=1$, and $x=(1,2,-1,3)$, both execution forms yield
$(1,2.5,0.25,3.125)$. With a nonzero $h_0$, Eq.~\ref{eq:convolution}
also needs the term $ca^t h_0$. Omitting it is a state-contract error.

\Plate{03-kernel}{A causal kernel for $a=0.5$ and unit input/readout gain. The lower-triangular coefficients depend only on lag. The diagram shows why future inputs are excluded and why fixed coefficients allow a shared convolution kernel.}{fig:kernel}

A diagonal $N$-coordinate state produces a sum of exponential kernels,
$K_\ell=\sum_n c_n g_n a_n^\ell$. Those coordinates may encode different
decay rates or, in other parameterizations, oscillations. Their sum
can represent richer filters than any individual coordinate.
Efficient kernel computation, parameter initialization, and numerical
conditioning are separate questions from this algebraic equivalence.

\section{Make updates selective without making them state dependent}

Suppose $\Delta_t$ is computed from the current input:
\begin{equation}
 \Delta_t=\operatorname{softplus}(W_\Delta x_t+b_\Delta),
 \quad a_t=e^{-\lambda\Delta_t},
 \quad b_t=\frac{1-a_t}{\lambda}u_t.
\end{equation}
The update remains affine in $h_{t-1}$, but its coefficients change with
the input. Small $\Delta_t$ retains state and weakens the write.
Large $\Delta_t$ forgets the old state and moves toward $u_t/\lambda$.
It does not generally reset the state to zero.

\Listing{selected_coefficients}

Our small example reuses a projection for the drive and interval selection
to keep the calculation inspectable. That choice is a teaching simplification,
not the parameterization of a published architecture. Each coefficient
depends on $x_t$ and fixed parameters, not on the unknown previous recurrent
state. All coefficients can therefore be formed before scanning this layer.
They can be nonlinear functions of the input without breaking that property.

Why does a fixed convolution no longer suffice? Expanding the variable
recurrence gives
\begin{equation}
 h_t=\left(\prod_{k=1}^t a_k\right)h_0+
      \sum_{j=1}^t\left(\prod_{k=j+1}^t a_k\right)b_j.
 \label{eq:variable}
\end{equation}
The contribution of input $j$ depends on intervening inputs through the
$a_k$. There is generally no single lag-only kernel shared across
sequences. Even the one-step map
$y(x)=(1-e^{-\operatorname{softplus}(x)})x$ fails
$y(2x)=2y(x)$, so it is not a linear operator on the input.

Mamba's selective SSM makes certain state-space parameters input dependent
and uses scan-based execution; its architecture and hardware-aware
implementation add further mechanisms \cite{mamba}. We use the principle
to motivate an original toy, not to claim equivalent accuracy or speed.

Contrast this with Chapter 6's GRU: its gates and candidate depend
nonlinearly on $h_{t-1}$. The GRU's next affine-looking pair cannot
generally be precomputed independently of earlier states. Function
composition is always associative abstractly, but composing arbitrary
nonlinear functions need not yield a small, cheaply evaluable summary.
The useful property here is closure in a fixed-size family of affine maps.

\section{Compose intervals, preserving chronological order}

Let an earlier interval map $h$ to $ah+b$, and a later interval map
its input to $ch+d$. The composition is
\begin{equation}
 c(ah+b)+d=(ca)h+(cb+d).
\end{equation}
Define the chronological binary operation
\begin{equation}
 (a,b)\star(c,d)=(ca,cb+d).
 \label{eq:compose}
\end{equation}
The first argument occurs earlier in time. This convention is important:
some software APIs reverse the argument order while expressing the
same functional composition.

\Listing{compose}

For three pairs $(a,b),(c,d),(e,f)$, either parenthesization gives
\begin{equation}
 (eca,\ ecb+ed+f).
\end{equation}
Thus $\star$ is associative over exact real arithmetic. Its identity
is $(1,0)$. It is not commutative: $(0.5,1)\star(0.25,2)$ gives
$(0.125,2.25)$, while reversing them gives $(0.125,2)$.
Changing the grouping is legal; changing the order changes the sequence.

For vectors with diagonal transitions, the products are elementwise.
For dense matrices the same argument uses ordered matrix multiplication.
Associativity survives, but composing dense transitions costs much more
than composing diagonal ones. A proof of parallelizability does not
remove the cost of the operator used inside the scan.

\section{Compute every prefix with a tree}

An inclusive \Term{prefix scan} returns each prefix composition:
$p_1$, $p_1\star p_2$, and so on. Blelloch develops scans and reductions
of recurrences to associative operations \cite{blelloch}.
Our implementation first composes adjacent pairs, recursively scans those
pair summaries, then fills the intermediate prefixes.

For eight inputs, the recursive sizes are $8,4,2,1$. At the size-eight
level we form four pair summaries. The recursive result gives prefixes
ending at positions 2, 4, 6, and 8. Position 1 is already known;
positions 3, 5, and 7 each require one more composition. Nothing is
reordered, and no future coefficient participates in an earlier prefix.

\Plate{04-scan-tree}{A four-element inclusive scan. The reduction combines adjacent intervals; reconstruction supplies the missing prefix ending at element three. Labels identify the interval represented by each summary. The arithmetic is associative, but chronological order remains fixed.}{fig:scan}

\Listing{prefix_pairs}
\Listing{parallel_form}

If prefix $t$ is $(P_t,Q_t)$, applying it to the initial state produces
$P_t\odot h_0+Q_t$. Keeping $h_0$ outside the prefix recursion makes
nonzero initial-state handling explicit and supports chunk summaries.
It also exposes a common bug: returning only $Q_t$ accidentally assumes
zero initial state.

For $T\geq2$, our number of pair compositions satisfies
\begin{equation}
 W(T)=\lfloor T/2\rfloor+
      \lfloor(T-1)/2\rfloor+W(\lfloor T/2\rfloor),\quad W(1)=0.
\end{equation}
Hence $W(T)<2T$. For $T=8$, the count is 11, not the seven
compositions needed merely to reduce to one final summary. A scan
returns all prefixes, so it solves a larger output problem.

The dependency depth is $O(\log T)$ when independent nodes at each level
are executed concurrently. \emph{This Python reference does not do that.}
Its loops execute on the host; the tensor operations expose the dependency
structure but are not a fused accelerator kernel. The algorithm has
linear arithmetic work for fixed state width, but these measurements
do not establish wall-clock speedups.

\section{Check gradients and finite-precision equivalence}

Changing parenthesization can change floating-point rounding.
Associativity is a real-arithmetic identity, not a promise of bitwise
equality between execution orders. We compare CPU float64 outputs with
tolerances, including empty, odd-length, and non-power-of-two sequences.
Every prefix is checked, not only the last one.

The tests compare gradients of a squared-output loss with respect to
all $a_t$, $b_t$, and $h_0$. They also apply finite-difference gradient
checking to the scan and to the input-dependent coefficient path.
This catches a hidden detachment that could preserve the forward result
while changing the learned model.

For fixed coefficients, the direct state sensitivity is
\begin{equation}
 \frac{\partial h_t}{\partial h_0}=\prod_{k=1}^t a_k.
\end{equation}
When coefficients depend on inputs, input gradients additionally traverse
the functions producing $a_k$ and $b_k$. Treating those coefficients as
constants during differentiation would train a different computation.
Our out-of-place construction preserves those paths.

Positive decay and nonnegative intervals keep $0<a_t\leq1$ in exact
arithmetic. This bounds the homogeneous retention path but does not
bound every output or gradient: the drive, readout, projections, and
their derivatives also matter. Extremely long products can underflow.
Testing modest float64 examples establishes a reference, not universal
mixed-precision stability.

\section{Make padding, reset, and carry algebraic}

For finite inputs, padding that must not advance state is the identity
pair $(1,0)$. It preserves the incoming state exactly at the mathematical
level. Feeding a zero token through learned projections need not produce
that pair. Loss masking remains separate: padded outputs must not become
valid targets merely because their state was held.

A record boundary that resets the incoming state to zero \emph{before}
the current input's update replaces $(a_t,b_t)$ by $(0,b_t)$.
Then $h_t=b_t$, with no contribution from the preceding record.
Resetting after the update would instead discard the current input.
At a chunk boundary within one record, carry the final state; do not
insert either padding or reset.

\Plate{05-boundaries}{Three different boundary contracts expressed as affine maps. Padding is identity, record reset removes the incoming-state term, and chunk carry preserves the actual boundary state. These operations can have identical tensor shapes while representing different computations.}{fig:boundaries}

Full and chunked execution match when coefficients, token positions,
model parameters, and initial-state policy match. A coefficient network
with a local causal convolution would also need its own boundary buffer;
carrying only the SSM state would then be insufficient. A change in
parameters between chunks creates the same stale-state issue discussed
in Chapter 6. These details become explicit runtime contracts later.

\section{Place the mechanism in the broader architecture}

S4, S5, Mamba, and RWKV should not be collapsed into one model merely
because they use recurrence. S4 develops structured state-space sequence
models. S5 uses a multi-input, multi-output state-space layer and parallel
scan. The original Mamba combines selective state spaces with an
architecture and a hardware-aware execution strategy. The original
RWKV paper presents a recurrent language-model architecture with its own
time-mixing and channel-mixing construction \cite{s4,s5,mamba,rwkv}.
These are high-level distinctions; this chapter does not reproduce
their benchmark results or implement their complete blocks.

For Evolutor, the reusable lesson is a legality test:
can an interval's effect be represented in a fixed-size family whose
composition is associative and affordable? If yes, a compiler may have
sequential, chunked, and scan execution choices. If not, a label such as
state model does not make those transformations legal.

DOGMA remains the non-Transformer DNA-native research direction;
Hermon DNA remains the Transformer-based direction. Neither identity
is established by giving a generic diagonal recurrence a biological name.
Claims about genomic usefulness require tasks, appropriate data splits,
baselines, and mechanism-specific interventions.

Run the reference and tests:
\begin{Verbatim}[fontsize=\small]
python3 drafts/ch07/reference.py
python3 -m pytest -q tests/test_ch07_reference.py
\end{Verbatim}
Chapter 8 changes the question from how to compress history into a
fixed-dimensional state to how to address stored content through attention.
Chapter 26 will revisit scan parity at model scale, with training,
precision, and runtime boundaries included.

\section{Exercises with worked solutions}

\subsection*{1. Trace an identity}
Apply $(0.5,1)$ and then $(1,0)$ to $h_0=2$.
\textbf{Solution:} the first state is 2; the identity leaves it at 2.
An identity step is not a reset to zero.

\subsection*{2. Reverse two updates}
Compare $(0.5,1)\star(0.25,2)$ with the reversed order.
\textbf{Solution:} the multipliers both equal 0.125, but the writes
are 2.25 and 2. Associativity permits regrouping, not permutation.

\subsection*{3. Audit Euler stability}
For $\lambda=2$ and $\Delta=1$, compare Euler and exact retention.
\textbf{Solution:} Euler gives $-1$, while the exact multiplier is
$e^{-2}\approx0.1353$. Euler does not decay in magnitude in this case.

\subsection*{4. Recover a limiting gain}
Show that $(1-e^{-\lambda\Delta})/\lambda\to\Delta$ as
$\lambda\Delta\to0$ with fixed positive $\lambda$.
\textbf{Solution:} expand the exponential:
$1-e^{-\lambda\Delta}=\lambda\Delta+O(\Delta^2)$.
The implementation uses \texttt{expm1} to preserve this small difference.

\subsection*{5. Restore a nonzero initial state}
Which term is missing from a zero-state convolution when $h_0\ne0$?
\textbf{Solution:} add $ca^t h_0$ to output $y_t$ for fixed coefficients.
The prefix form handles this as $P_th_0+Q_t$ before readout.

\subsection*{6. Find the hidden nonlinearity}
Why can a model be affine in state but nonlinear in inputs?
\textbf{Solution:} the coefficient functions may depend nonlinearly on
the input. Holding those coefficients fixed is a conditional linear
calculation, not a claim that the entire input-output map is linear.

\subsection*{7. Count a scan}
How many compositions does our size-eight scan perform?
\textbf{Solution:} $(4+3)+(2+1)+(1+0)=11$.
A final-only reduction needs seven, but does not return every prefix.

\subsection*{8. Check an odd length}
At length five, where does the final prefix come from?
\textbf{Solution:} scan the two complete adjacent pairs, then compose
the prefix covering positions 1--4 with the original fifth pair.
No fake padded input is needed.

\subsection*{9. Detect a gradient bug}
The scan matches sequential outputs but its coefficient network has
zero gradients. What should be tested?
\textbf{Solution:} compare parameter and input gradients and check for
detachment or conversion to Python numbers between coefficient creation
and scanning. Forward parity alone cannot expose all such errors.

\subsection*{10. Encode a reset}
Which pair resets to zero before applying the current write $b_t$?
\textbf{Solution:} $(0,b_t)$. The new state is $b_t$.
Using $(0,0)$ would discard the current write too.

\subsection*{11. Challenge a speed claim}
Does logarithmic dependency depth prove this Python code is faster?
\textbf{Solution:} no. It specifies a parallel schedule's potential
depth, not available processors, kernel overhead, memory traffic, or
measured latency. Report both theoretical work and actual execution.

\subsection*{12. Design a fair memory comparison}
How would you compare a selective recurrence with attention on DNA?
\textbf{Worked rubric:} align tokenization, targets, splits, training
budgets, and evaluation distances. Test state reset and causality;
separate synthetic retention from learned accuracy. Measure inference
memory and latency under declared batch and sequence lengths. Include
negative controls for selection and readout rather than attributing
every improvement to recurrence.
